\section{Introduction}\label{sec:intro}

Stablecoins have become the settlement asset of the crypto economy and, increasingly, a dollar substitute for people
outside the United States. They have also become the preferred medium of the industrialised online fraud that grew
out of Southeast Asia's border casinos and special economic zones. Regional and international bodies describe a
system in which victims' bank transfers are converted, through networks of rented accounts, into USDT on the Tron
blockchain and moved through informal ``guarantee'' marketplaces and underground bankers before being integrated into
real assets \citep{unodc2024casinos,unodc2025inflection,fatf2023cyberfraud}. Thailand sits at the edge of this system.
It hosts a licensed, identity-verified crypto exchange industry and a large retail investor base; it shares long land
borders with the compounds of Myawaddy, the casinos of Poipet and the special economic zones of the Golden Triangle;
and its banks, police and regulators have spent the past four years tightening controls on so-called mule accounts.

This paper asks three questions that a Thai regulator, an exchange compliance officer or a researcher would want
answered. First, is Thai stablecoin behaviour unusual---relative to other crypto-assets traded on the same platforms
and relative to stablecoin use elsewhere? Second, is it plausible that criminal networks, including the ``call
centre'' fraud gangs that target Thai victims, launder proceeds through USDT, and through Thai rails in particular?
Third, and most important, how could either claim be tested in a way that separates criminal flows from the many
legitimate reasons to hold and move digital dollars?

The third question disciplines the first two. Unusual behaviour is not evidence of crime. Thai investors cannot trade
crypto derivatives on licensed domestic venues, so buying USDT and withdrawing it to an offshore exchange is the
ordinary route to leverage. Millions of migrant workers from Myanmar, Cambodia and Laos live in Thailand and send money
home, and two of those countries have recently experienced the high inflation and capital controls that make dollar
stablecoins attractive. Honest peer-to-peer dealers receive baht from many counterparties and withdraw quickly, which
is exactly what a mule account does. A credible study must therefore name each legitimate motive, state the
fingerprint it leaves in data, and show that the residual pattern attributed to laundering survives once those
fingerprints are removed.

\paragraph{What this paper does}
I make four contributions. The first is \emph{descriptive and uses only public data}. Using the Coin Metrics community
files through 23 May 2026, I document that USDT on Tron behaves as a retail payment rail that is economically distinct
from USDT on Ethereum (\cref{sec:rails}). Tron carries 47\% of the USDT supply on the chains covered by the data, up
from 19\% at the end of 2019. It hosts 74.5 million addresses with a positive balance---5.3 times as many as USDT on
Ethereum---with a mean balance of about US\$1,200 against roughly US\$7,000, and it records 9.3 times as many transfers
per dollar of supply. When Circle wound down USDC on Tron in 2024, USDC supply there fell from a peak of US\$3.5
billion to under US\$30 million, leaving the cheapest high-liquidity dollar rail in the region to a single issuer. In
the U.S.\ Treasury's sanctions list, 95\% of the addresses tagged USDT, TRX or USDC are Tron addresses. I also show that
Thailand's macroeconomic fundamentals give its residents little reason to hold dollars on-chain---its consumer-price
inflation averaged 1.7\% a year over 2019--2023, the fourth lowest of 23 comparators---whereas Laos (14.9\%) and Myanmar
have strong legitimate demand (\cref{sec:macro}).

The second contribution is \emph{a negative result with a methodological lesson}. Event studies on chain-wide activity
are a tempting way to evaluate crackdowns because the data are free. I estimate abnormal Tron USDT activity around
three dated shocks---Thailand's February 2025 power cut to border scam hubs, its June 2025 block on unlicensed offshore
exchanges, and the October 2025 U.S.\ and U.K.\ sanctions on Prince Group---and calibrate inference against 149 placebo
dates. Conventional bootstrap bands, which already account for serial correlation and trend uncertainty, reject the
null at 58\% of random dates. Calibrated against the placebo distribution, none of the three shocks is distinguishable
from an ordinary week ($p$ between 0.69 and 1.00). \Cref{sec:displacement} rationalises this: in a channel-choice model
in which enforcement raises the cost of one rail, flows are displaced to substitutes and the aggregate is unchanged to
first order.

The third contribution is \emph{a set of estimators matched to the economics of structuring and pass-through}
(\cref{sec:framework}). The central one exploits the Bank of Thailand's face-verification rule, which applies to
transfers of 50,000 baht or more. Honest users also dislike the scan and bunch below the threshold, so bunching by
itself proves nothing. Under an assumption that hassle avoidance is the same across the assets a deposit funds, the
\emph{difference} in a slope-invariant bunching statistic between USDT-funding and other deposits identifies the volume
of structured deposits (\cref{prop:bunching}). I complement it with dwell-time mixtures, a direction-aware measure of
graph exposure to flagged addresses, and synthetic control for policy shocks, and I prove that undirected hop counts
mechanically flag exchange hubs (\cref{prop:direction}).

The fourth contribution is \emph{validation}. Because no public dataset links Thai exchange accounts to criminal
outcomes, I build an agent-based simulator of an exchange-year with five honest populations (domestic traders,
offshore bridge traders including arbitrageurs, cross-border remitters, holders and honest peer-to-peer merchants) and
mule accounts whose tradecraft varies, including camouflaged mules on rented, aged accounts (\cref{sec:sim}). Against
planted ground truth, the undirected two-hop rule attains a precision of 2--6\% among its top-ranked accounts; speed
alone attains 6--21\%; a direction-aware exposure measure rises from 2\% to 86\% as prevalence grows; and multi-signal
scoring reaches 78--92\%. The bunching-difference test rejects in 7.5\% of null simulations at a nominal 5\% level and
in 88\% of simulations with 2\% mule accounts. Synthetic control has a 2.5\% false-rejection rate and 82\% power for a
0.10 log-point effect. These magnitudes are properties of the simulator; the ranking of methods and their failure modes
are the transferable result.

\paragraph{What this paper does not claim}
The paper does not estimate how much Thai stablecoin volume is criminal. Doing so requires account-level data that
exchanges and the Anti-Money Laundering Office hold and that the Personal Data Protection Act rightly restricts.
Instead, \cref{sec:plan} sets out a pre-registered plan---six hypotheses, the data tier each needs, the estimator and
the decision rule---so that the question can be answered as soon as such data are available, without specification
search. All code is open, every figure is regenerated from raw files whose checksums are recorded, and every
simulated figure is marked as such.

\paragraph{Roadmap}
\Cref{sec:background} describes the Thai regime and the regional scam economy. \Cref{sec:literature} relates the paper
to work on stablecoins, illicit crypto finance and blockchain forensics. \Cref{sec:framework} presents the conceptual
framework, the model and the propositions. \Cref{sec:data} describes the data. \Cref{sec:evidence} reports the
evidence from public data and \cref{sec:sim} the simulation study. \Cref{sec:plan} contains the pre-registered plan,
policy implications, limitations and conclusions. Proofs, simulator parameters and replication instructions are in the
appendix.
